\clearpage
\begin{summarybox}
\textbf{\textcolor{CSummary}{一句话核心}}

\textbf{弦长 $=$ 坐标投影差 $\times$ 由直线方向决定的放大系数；
需要求交点时，优先“联立 $\rightarrow$ 根差 $\rightarrow$ 弦长”，不用先解两个点。}

\medskip
\textbf{\textcolor{CSummary}{完整认知链（初中 $\rightarrow$ 高中）}}
\[
 \underbrace{|AB|=2\sqrt{r^2-d^2}}_{\text{初中：圆的垂径定理}}
 \quad\xrightarrow{\text{保留勾股，换直角边}}\quad
 \underbrace{|AB|=\sqrt{(\Delta x)^2+(\Delta y)^2}}_{\text{所有平面线段}}
\]
再使用 $\Delta y=k\Delta x$，便得到
$|AB|=\sqrt{1+k^2}|\Delta x|$。
\textbf{一句话：}不是从“圆弦公式”死记出新公式，
而是从“圆弦背后的勾股定理”迁移出通用公式。

\medskip
\textbf{\textcolor{CSummary}{最小公式组}}
\[
 \boxed{\begin{aligned}
 y=kx+b:&\quad |AB|=\sqrt{1+k^2}\,|\Delta x|;\\
 x=ty+d:&\quad |AB|=\sqrt{1+t^2}\,|\Delta y|;\\
 \alpha z^2+\beta z+\gamma=0:&\quad
 |z_1-z_2|=\frac{\sqrt{\beta^2-4\alpha\gamma}}{|\alpha|}.
 \end{aligned}}
\]
第三式\textbf{要求} $\alpha\ne0$、$\Delta>0$。
水平线优先横差；竖直线优先纵差。

\medskip
\textbf{\textcolor{CSummary}{三种证明，一次掌握}}
\begin{itemize}[leftmargin=*,itemsep=0.25em]
 \item \textbf{最快：}斜率 $\Delta y=k\Delta x$，代入勾股，一步提取 $|\Delta x|$。
 \item \textbf{门槛最低：}看直角三角形，两直角边比例为 $1:|k|$，斜边比例为 $\sqrt{1+k^2}$。
 \item \textbf{常规：}坐标距离公式代入直线方程，并强调平方开方取绝对值。
\end{itemize}

\medskip
\textbf{\textcolor{CSummary}{考场速写流程}}

\textbf{判断两交点 $\rightarrow$ 看消 $x$ 还是消 $y$ 更简单
$\rightarrow$ 检查二次项及 $\Delta>0$
$\rightarrow$ 由韦达或 $\sqrt\Delta/|\alpha|$ 求根差
$\rightarrow$ 乘方向系数。}

\medskip
\textbf{\textcolor{CSummary}{两个最小反例}}
\begin{itemize}[leftmargin=*,itemsep=0.2em]
 \item $k=0$ 却使用含 $1/k^2$ 的纵坐标公式：\textbf{无定义}。
 \item $x^2-y^2=1$ 与 $y=x+1$：消元退化一次，\textbf{不能当作两交点弦}。
\end{itemize}

\textbf{\textcolor{CSummary}{迁移原则}}
同一套投影公式适用于圆、椭圆、双曲线、抛物线，
实际上适用于同一直线上的\textbf{任意两点}。
\end{summarybox}
